
\subsection*{4.5}
On
appelle équivalence numérique la relation d'équivalence définie par le produit d'intersection \(\langle\, ,\,\rangle\) :

\textbf{Définition 4.5.1.} \uline{On dit que deux éléments} \(x\) \uline{et} \(y\) \uline{de} \(\Gr^i(X)\) \((i\in Z)\) (resp. \uline{de} \(\Gr_i(X)\)\uline{) sont numériquement équivalents, s'ils définissent le même homomorphisme} \(\Gr_i(X)\longrightarrow\mathbb Z\) (resp. \(\Gr^i(X)\longrightarrow\mathbb Z\)\uline{).
On dit que deux} \(\underline{O}_X\)\uline{-modules inversibles} \(\underline{L}_1\) et \(\underline{L}_2\) \uline{sont numériquement équivalents, si leurs premières classes de Chern} \(x,y\in\Gr^1(X)\) \uline{sont numériquement équivalentes.}

\textbf{Proposition 4.5.2.} \uline{Un} \(\underline{O}_X\)\uline{-module inversible} \(\underline{L}\) \uline{est numériquement équivalent à zéro si et seulement si}
\begin{equation}\tag{4.5.3}
 \chi(\underline{L}^n\otimes_{\underline{O}_X}\underline{O}_Z)=\chi(\underline{O}_Z)
\end{equation}
\uline{pour tout} \(n\in\mathbb Z\) \uline{et pour toute courbe intègre fermée} \(Z\subset X\).

Cela résulte de 1.1.2 et 4.3.1.

On désigne par \(\Pic^N(X)\) le sous-groupe de \(\Pic(X)\) composé des classes de \(\underline{O}_X\)-modules inversibles numériquement équivalents à zéro. De 2.3.4 on tire

\textbf{Corollaire 4.5.3.} \(\Pic^\tau(X)\subset \Pic^N(X)\).

On va démontrer dans Exp XIII que l'on a même \(\Pic^N(X)=\Pic^\tau(X)\).

\section*{5. L'isomorphisme \(\Pic(X)\cong\Gr^1(X)\)}

\subsection*{5.1.}
Soit \(X\) un topos annelé quelconque. Alors on peut définir l'anneau \(K^{\bullet}(X)\) et la filtration \((\Filt^i(X))_{i\in\mathbb Z}\) exactement comme pour les schémas. De même on définit le groupe \(\Pic(X)\) comme le groupe des classes d'isomorphie des \(\underline{O}_X\)-modules inversibles. On désigne par
